\section{Visão Geral e Princípios Arquiteturais}

O Núcleo de Situação de Saúde (NSS) é organizado como uma aplicação web apoiada por um backend Java/Spring Boot, um pipeline de dados em Python e um banco de dados PostgreSQL. O objetivo desta documentação é registrar as decisões que permitem que os três componentes evoluam de forma independente, mas permaneçam integrados como um único sistema.

\subsection{Objetivo da arquitetura}

A arquitetura separa duas responsabilidades que possuem características diferentes:

\begin{itemize}
    \item \textbf{consulta da aplicação}: realizada pelo frontend por meio da API Java, com leitura dos dados consolidados no PostgreSQL;
    \item \textbf{ingestão e atualização dos dados}: realizada pelo pipeline Python, que obtém dados do SINAN por meio do PySUS, processa os dados e atualiza o PostgreSQL.
\end{itemize}

Assim, o Python não faz parte do caminho síncrono de uma consulta feita pelo usuário. O usuário consulta dados que já foram processados e persistidos.

\begin{regra}[Princípio de separação de responsabilidades]
O frontend não acessa diretamente o PostgreSQL nem o pipeline Python. A API Java é o ponto de entrada da aplicação. O pipeline Python é responsável pela preparação e atualização dos dados.
\end{regra}

\subsection{Princípios}

\begin{enumerate}
    \item Cada componente possui uma responsabilidade técnica bem definida.
    \item Os três componentes de aplicação permanecem em repositórios independentes.
    \item O Docker fornece ambientes reproduzíveis para desenvolvimento e integração.
    \item O PostgreSQL é o armazenamento compartilhado dos dados utilizados pela aplicação.
    \item DataFrames e arquivos Parquet permanecem no pipeline para processamento, rastreabilidade, inspeção e debugging.
    \item O pipeline deve permitir identificar em qual etapa uma inconsistência foi introduzida.
    \item Contratos de dados e mudanças de schema devem ser tratados de forma versionada.
    \item Segurança, configuração e credenciais não devem ser incorporadas ao código-fonte.
\end{enumerate}

\begin{explicacao}[Decisão arquitetural]
A separação em três repositórios não implica necessariamente três aplicações isoladas. Os componentes podem ser desenvolvidos separadamente e integrados em um único ambiente por meio de Docker Compose, mantendo uma rede Docker compartilhada e containers distintos.
\end{explicacao}
